\newpage
\phantomsection
\label{prf:euclidean-metric}
\LRAProofFor{cor:euclidean-metric}

\begin{remark*}[Return]
\hyperref[cor:euclidean-metric]{Return to Corollary}
\end{remark*}

\begin{corollary*}[Euclidean Metric]
For \(x,y\in\mathbb{R}^n\), define
\[
  d_2(x,y)=\sqrt{\sum_{i=1}^{n}|x_i-y_i|^2}.
\]
Then \((\mathbb{R}^n,d_2)\) is a metric space.
\end{corollary*}

\begin{proof}[Professional Standard Proof]
\LRAProofBodyStart
The function \(d_2\) is the Euclidean distance on \(\mathbb{R}^n\).  For all
\(x,y\in\mathbb{R}^n\), \(d_2(x,y)\geq 0\), and \(d_2(x,y)=0\) if and only if
each coordinate difference \(x_i-y_i\) is \(0\), equivalently \(x=y\).
Symmetry follows from \(|x_i-y_i|=|y_i-x_i|\) for every coordinate.

For the triangle inequality, write
\[
  x-y=(x-z)+(z-y).
\]
Applying \hyperref[cor:minkowski-inequality-rn]{Minkowski's inequality} gives
\[
  d_2(x,y)\leq d_2(x,z)+d_2(z,y).
\]
Thus \(d_2\) satisfies the metric axioms, so \((\mathbb{R}^n,d_2)\) is a
metric space.
\end{proof}

\begin{proof}[Detailed Learning Proof]
\LRAProofBodyStart
The formula for \(d_2\) is the length of the vector \(x-y\).  A length is
nonnegative.  It is zero exactly when every coordinate of \(x-y\) is zero,
which means \(x=y\).  It is symmetric because replacing \(x-y\) by \(y-x\)
only changes signs before absolute values and squares are taken.

The only structural point is the triangle inequality.  The displacement from
\(y\) to \(x\) can be split through \(z\):
\[
  x-y=(x-z)+(z-y).
\]
Minkowski's inequality says that the length of a sum is at most the sum of
the lengths.  Therefore the direct distance \(d_2(x,y)\) is at most
\(d_2(x,z)+d_2(z,y)\).
\end{proof}

\begin{remark*}[Proof structure]
The proof checks the easy metric axioms coordinatewise and uses Minkowski's
inequality for the triangle inequality.
\end{remark*}

\begin{dependencies}
\begin{itemize}
  \item \hyperref[def:metric-space]{Definition: Metric Space}
  \item \hyperref[def:euclidean-distance-rn]{Definition: Euclidean Distance on \(\mathbb{R}^n\)}
  \item \hyperref[cor:minkowski-inequality-rn]{Corollary: Minkowski's Inequality}
\end{itemize}
\end{dependencies}

\clearpage
